\documentclass[letterpaper]{article} % DO NOT CHANGE THIS
\usepackage[submission]{aaai2027}  % DO NOT CHANGE THIS
\usepackage{times}  % DO NOT CHANGE THIS
\usepackage{helvet} % DO NOT CHANGE THIS
\usepackage{courier}  % DO NOT CHANGE THIS
\usepackage[hyphens]{url}  % DO NOT CHANGE THIS
\usepackage{graphicx} % DO NOT CHANGE THIS
\urlstyle{rm} % DO NOT CHANGE THIS
\def\UrlFont{\rm}  % DO NOT CHANGE THIS
\usepackage{natbib}  % DO NOT CHANGE THIS AND DO NOT ADD ANY OPTIONS TO IT
\usepackage{caption} % DO NOT CHANGE THIS AND DO NOT ADD ANY OPTIONS TO IT
\frenchspacing  % DO NOT CHANGE THIS
\setlength{\pdfpagewidth}{8.5in}  % DO NOT CHANGE THIS
\setlength{\pdfpageheight}{11in}  % DO NOT CHANGE THIS

% Recommended for better-looking tables.
\usepackage{booktabs}

\pdfinfo{
/TemplateVersion (2027.1)
}

\setcounter{secnumdepth}{0}

\title{Ink Density Saturation in Duan Inkstones}
% An anonymous submission names no authors and leaves affiliations empty,
% but the \affiliations block itself is required: \maketitle expands
% \aaai@affiliations unconditionally.
\author{
    Anonymous Submission
}
\affiliations{
}

\begin{document}

\maketitle

\begin{abstract}
A minimal document in the AAAI 2027 submission format. The content
mirrors \texttt{latex-rendering-check.tex}: ink density saturates as a
function of grind time, and the stone --- not the stick --- sets the
ceiling.
\end{abstract}

\section{Model}

Let $D(t)$ denote ink density after $t$ minutes of grinding. Density
follows a saturating exponential,
\begin{equation}
  \label{eq:density}
  D(t) = D_{\max}\left(1 - e^{-t/\tau}\right),
\end{equation}
where $\tau$ is the stone's time constant and $D_{\max}$ its ceiling.
The relationship was established empirically by \citet{tanaka2019}.

\section{Results}

Table~\ref{tab:stones} lists time constants for three stones. The
grinding surface, not the ink stick, dominates.

\begin{table}[t]
\centering
\begin{tabular}{lrr}
\toprule
Stone & $\tau$ (min) & $D_{\max}$ \\
\midrule
Duan (Laokeng)  & 4.1 & 0.94 \\
Duan (Majiakeng) & 6.8 & 0.87 \\
She             & 9.2 & 0.81 \\
\bottomrule
\end{tabular}
\caption{Saturation parameters by stone.}
\label{tab:stones}
\end{table}

\section{Checklist}

Opening this file and pressing preview should produce a two-column PDF
with the AAAI running head, a numbered equation, a booktabs table, and a
resolved citation to \citet{tanaka2019} in the references.

% aaai2027.sty sets the bibliographystyle itself; declaring it again makes
% bibtex fail with "Illegal, another \bibstyle command".
\bibliography{references}

\end{document}
